%
% Halaman Abstrak
%
% @author  Andreas Febrian @version 2.00 @edit by Ichlasul Affan
%

\chapter*{Abstrak}
\singlespacing

\vspace*{0.2cm}

\noindent \begin{tabular}{l l p{10cm}} Nama & : & \penulis \\
	Program Studi      & : & \program \\
	Judul              & : & \judul   \\
\end{tabular} \\

\vspace*{0.5cm}

\noindent Peraturan perundang-undangan pada umumnya tersedia dalam bentuk berkas PDF yang bersifat
tidak \textit{machine-readable}. \textit{Knowledge graph} adalah struktur data berbentuk
\textit{graph} yang menggambarkan entitas dunia nyata beserta keterkaitannya. Pada penelitian ini
berfokus pada pengembangan \textit{framework} untuk mengonversi peraturan perundang-undangan di
Indonesia dari dokumen PDF menjadi \textit{knowledge graph} dan memberikan contoh \textit{use case}
dari \textit{knowledge graph} tersebut. \textit{Knowledge graph} peraturan perundang-undangan
tersebut mengandung data-data berupa metadata peraturan, struktur peraturan, teks peraturan, serta
relasi antara peraturan seperti amandemen dan rujukan. 

// TODO: Sblm kalimat ini, jelasin jg scara singkat statistik berapa jmlh peraturan yg bs kita
ekstraksi dan brp jmlh triplesnya. Lalu jelasin jg klo Aab dlm angka statistik tsb, Aab salah
satunya bs menangani UU Cipta Kerja yang bersifat kompleks dan berukuran besar. Ini penting karena
menunjukkan pencapaian dari skripsi Aab.

Peraturan perundang-undangan dalam bentuk
\textit{knowledge graph} dapat dimanfaatkan dalam berbagai penerapan yang menggunakan komputer
seperti \textit{question answering}, \textit{reasoning}, dan visualisasi data.

// TODO: Ini terdengar agak pasif, mgkn yg lbh aktif, misal: Penelitian ini juga mendemonstrasikan
pemanfaatan knowledge graph tersebut untuk pastikan pemanfaat yg ditulis di sini ada di Bab use
case nantinya, jd biar promise yg ada di abstrak, terjawab di body skripsi Aab


\vspace*{0.2cm}

\noindent \bo{Kata kunci:} \\ \f{Knowledge Graph}, Peraturan perundang-undangan, SPARQL, \f{Resource
Definition Framework}

\onehalfspacing
\newpage
